% ==========================================================================
%  Guía de Estudio · M03 · Potencial Eléctrico y Superficies Equipotenciales
%  Electro10.Easy — Universidad Autónoma de Bucaramanga
%  Autor: Prof. Dudbil Olvasada Pabón Riaño · 4 páginas exactas
%  Compilar:  pdflatex M03_Guia.tex   (desde .design/fase5/M03/)
% ==========================================================================
\documentclass[11pt]{article}
\usepackage[utf8]{inputenc}
\usepackage[T1]{fontenc}
\usepackage[spanish,es-noshorthands]{babel}
\usepackage[a4paper,margin=1.7cm]{geometry}
\usepackage{e10}
\usepackage{enumitem}
\usepackage[hidelinks]{hyperref}
\usepackage{tabularx}
\usepackage{array}

\setlength{\parindent}{0pt}
\setlength{\parskip}{0.5em}

\begin{document}

% ═══════════════════════════ PÁGINA 1 ═══════════════════════════
\begin{center}
  {\color{e10primary}\rule{\linewidth}{3.5pt}}\\[1.2em]
  {\LARGE\bfseries Guía de Estudio · Módulo 03}\\[0.25em]
  {\huge\bfseries Potencial Eléctrico y Superficies Equipotenciales}\\[0.6em]
  {\color{e10muted}Electro10.Easy · Cátedra de Electromagnetismo · Prof. Dudbil Olvasada Pabón Riaño}\\[0.3em]
  {\color{e10muted}\small Rigor: Sears--Zemansky, \emph{Física Universitaria} Vol.~2, Cap.~23 · $\Delta V=-\int\vec E\cdot d\vec\ell,\ \ \vec E=-\vec\nabla V$}
  \\[1em]
  {\color{e10primary}\rule{\linewidth}{1pt}}
\end{center}

\esection{Ideas clave}
\begin{itemize}[leftmargin=1.4em,itemsep=2pt,topsep=2pt]
  \item \textbf{El potencial es energía por unidad de carga:} $V=U/q_0$, en volt ($1\ \mathrm{V}=1\ \mathrm{J/C}$). Es un \textbf{escalar}.
  \item \textbf{El voltaje es una altura en el mapa:} $\vec E$ apunta cuesta abajo, de mayor a menor $V$. Una carga positiva ``cae'' hacia el potencial menor; la negativa, al revés.
  \item \textbf{Superficies equipotenciales} = curvas de nivel de $V$. $\vec E$ es \textbf{perpendicular} a ellas y mover carga sobre una no cuesta trabajo.
  \item \textbf{$V=0$ no implica $\vec E=0$:} en la mediatriz de un dipolo $V=0$ pero el campo es no nulo.
  \item \textbf{El electrón-volt es energía, no voltaje:} $1\ \mathrm{eV}=1{,}602\times10^{-19}\ \mathrm{J}$.
\end{itemize}

\esection{Organizador gráfico}
\renewcommand{\arraystretch}{1.25}
\begin{center}
\begin{tabular}{m{3.1cm} m{0.8cm} m{4.2cm} m{0.8cm} m{5.2cm}}
  \fcolorbox{e10primary}{e10primarysoft}{\parbox[c][1.4cm][c]{2.9cm}{\centering\small\textbf{Energía potencial}\\ $U=q_0V$}} &
  \centering\color{e10primary}\LARGE$\xrightarrow{\ \div q_0\ }$ &
  \fcolorbox{e10primary}{e10primarysoft}{\parbox[c][1.4cm][c]{4.0cm}{\centering\small\textbf{Potencial}\\ $V=\sum\dfrac{kq_i}{r_i}$}} &
  \centering\color{e10primary}\LARGE$\xrightarrow{\ -\vec\nabla\ }$ &
  \fcolorbox{e10primary}{e10primarysoft}{\parbox[c][1.4cm][c]{5.0cm}{\centering\small\textbf{Campo eléctrico}\\ $\vec E=-\vec\nabla V$}}\\[-0.3em]
  & & \color{e10primary}$\downarrow$ & & \color{e10primary}$\downarrow$\\[-0.3em]
  & & \fcolorbox{e10accent}{e10accentsoft}{\parbox[c][1.3cm][c]{4.0cm}{\centering\small\textbf{Equipotenciales}\\ $W=q_0\Delta V=0$}} & & \fcolorbox{e10accent}{e10accentsoft}{\parbox[c][1.3cm][c]{5.0cm}{\centering\small\textbf{Campo uniforme}\\ $\Delta V=Ed$}}\\[-0.3em]
  & & \color{e10primary}$\downarrow$ & & \\[-0.3em]
  & & \fcolorbox{e10ok}{e10oksoft}{\parbox[c][1.3cm][c]{4.0cm}{\centering\small\textbf{Electrón-volt}\\ $1\ \mathrm{eV}=1{,}602\times10^{-19}\ \mathrm{J}$}} & & \\
\end{tabular}
\end{center}

\etip{Analogía topográfica}{El voltaje es la altura de un mapa de montaña; el campo es la pendiente de la ladera. Una canica rueda cuesta abajo igual que una carga positiva ``cae'' hacia el potencial menor. La gravedad y la fuerza electrostática son ambas conservativas.}

\etip{¿Por qué el pájaro no se electrocuta?}{Posado en un cable de alta tensión, sus dos patas están al mismo potencial: $\Delta V\approx0$ y no circula corriente por su cuerpo. Lo que daña es la \emph{diferencia} de potencial, no el valor absoluto.}

\newpage
% ═══════════════════════════ PÁGINA 2 ═══════════════════════════
\esection{Fórmulas maestras}
\begin{center}
\renewcommand{\arraystretch}{1.5}
\begin{tabularx}{\linewidth}{>{\raggedright\arraybackslash}p{4.2cm} >{\centering\arraybackslash}X >{\raggedright\arraybackslash}p{5.6cm}}
\hline
\rowcolor{e10primarysoft} \textbf{Magnitud} & \textbf{Fórmula} & \textbf{Cuándo usarla} \\
\hline
Diferencia de potencial & $\Delta V=-\displaystyle\int_a^b \vec E\cdot d\vec\ell$ & Conocido el campo; integra a lo largo de la trayectoria. \\
Potencial de carga puntual & $V=\dfrac{kq}{r}$ & Una carga, con $V=0$ en el infinito. \\
Superposición & $V=\displaystyle\sum_i\dfrac{kq_i}{r_i}$ & Varias cargas puntuales (suma escalar). \\
Distribución continua & $V=\displaystyle\int\dfrac{k\,dq}{r}$ & Anillo, disco, varilla o cascarón. \\
Energía potencial & $U=q_0V=\displaystyle\sum_{i<j}\dfrac{kq_iq_j}{r_{ij}}$ & Energía de una carga o de un sistema. \\
Campo uniforme & $\Delta V=Ed$ & Entre placas paralelas con $\vec E$ constante. \\
Gradiente & $\vec E=-\vec\nabla V$ & Recuperar el campo a partir de $V$. \\
Electrón-volt & $1\ \mathrm{eV}=1{,}602\times10^{-19}\ \mathrm{J}$ & Energías de partículas (eV, keV, MeV). \\
\hline
\end{tabularx}
\end{center}

\esection{Cuatro errores comunes resueltos}

\eerror{Error 1 · ``Si $V=0$, entonces no hay campo.''}{
\textbf{La intuición correcta:} el potencial es un número que puede anularse por \emph{cancelación de signos}. En la mediatriz de un dipolo las dos contribuciones $kq/r$ y $k(-q)/r$ suman cero, pero el campo apunta de la carga positiva a la negativa. $V=0$ solo fija la referencia de altura; la pendiente ($\vec E=-\vec\nabla V$) puede ser distinta de cero ahí.
}

\eerror{Error 2 · ``Los potenciales se suman como vectores.''}{
\textbf{La intuición correcta:} el potencial es un \textbf{escalar}. $V=\sum_i kq_i/r_i$ es una suma de números reales con su signo, sin ángulos ni componentes. Esa es justamente su ventaja frente al campo: en el dipolo del ejemplo 23.4 de Sears--Zemansky, $V_a=1800+(-2700)=-900\ \mathrm{V}$, una suma algebraica directa.
}

\newpage
% ═══════════════════════════ PÁGINA 3 ═══════════════════════════
\eerror{Error 3 · ``El electrón-volt es una unidad de voltaje.''}{
\textbf{La intuición correcta:} el electrón-volt es una unidad de \textbf{energía}. $1\ \mathrm{eV}=1{,}602\times10^{-19}\ \mathrm{J}$ es la energía que gana una carga de magnitud $e$ al cruzar una diferencia de potencial de $1\ \mathrm{V}$. Un protón que atraviesa $7{,}5\ \mathrm{MV}$ gana $7{,}5\ \mathrm{MeV}$, no ``7,5 millones de electrones-volt de voltaje''.
}

\eerror{Error 4 · ``El potencial baja en contra del campo.''}{
\textbf{La intuición correcta:} de $\Delta V=-\int\vec E\cdot d\vec\ell$, al movernos \emph{a favor} de $\vec E$ el potencial \emph{disminuye}. El signo menos es la clave. Entre placas, $V$ cae de la placa positiva a la negativa, que es exactamente la dirección de $\vec E$.
}

\esection{Autoevaluación cuantitativa (5 preguntas)}
\begin{enumerate}[leftmargin=1.6em,itemsep=3pt]
  \item ¿Cuánto vale el potencial a $0{,}30\ \mathrm{m}$ de una carga de $+2{,}0\ \mathrm{nC}$? \hfill ($k=8{,}99\times10^9\ \mathrm{N\,m^2/C^2}$)
  \item Un electrón cruza una diferencia de potencial de $120\ \mathrm{V}$. ¿Cuánta energía gana, en eV y en joules?
  \item Un dipolo de $\pm12\ \mathrm{nC}$ separadas $10\ \mathrm{cm}$. ¿Cuánto vale $V$ en la mediatriz, a $13\ \mathrm{cm}$ de cada carga?
  \item Entre dos placas separadas $4{,}0\ \mathrm{mm}$ hay $\Delta V=80\ \mathrm{V}$. ¿Cuánto vale el campo $E$?
  \item Dado $V(x)=100-50x$ (en volt, con $x$ en metros), ¿cuánto vale $\vec E$ y hacia dónde apunta?
\end{enumerate}

\etip{Cómo usarla}{Resuelve primero en papel, luego contrasta con el solucionario de la página 4. Si fallas, vuelve al error común correspondiente antes de mirar la fórmula.}

\newpage
% ═══════════════════════════ PÁGINA 4 ═══════════════════════════
\esection{Solucionario numérico}

\esol{1.}{$V=kq/r=(8{,}99\times10^9)(2{,}0\times10^{-9})/0{,}30=59{,}9\ \mathrm{V}$.}
\esol{2.}{$\Delta K=|e|\,\Delta V=120\ \mathrm{eV}=120\times1{,}602\times10^{-19}=1{,}92\times10^{-17}\ \mathrm{J}$.}
\esol{3.}{$V=\dfrac{k(12\ \mathrm{nC})}{0{,}13}+\dfrac{k(-12\ \mathrm{nC})}{0{,}13}=830-830=0\ \mathrm{V}$. Ojo: $E\ne0$ en ese punto.}
\esol{4.}{$E=\Delta V/d=80/0{,}004=2{,}0\times10^4\ \mathrm{V/m}$.}
\esol{5.}{$E_x=-dV/dx=-(-50)=50\ \mathrm{V/m}$, apuntando en la dirección $+x$ (hacia donde $V$ disminuye).}

\esection{Errores detectados por telemetría (EquipotentialLab)}
\begin{itemize}[leftmargin=1.4em,itemsep=2pt]
  \item \texttt{potencial\_v\_vs\_u} — confundir $V$ con la energía $U$.
  \item \texttt{potencial\_r2} — usar $1/r^2$ en lugar de $1/r$.
  \item \texttt{potencial\_v0\_e0} — creer que $V=0$ implica $E=0$.
  \item \texttt{potencial\_equip\_perp} — dibujar equipotenciales no perpendiculares a $E$.
  \item \texttt{potencial\_ev\_vs\_v} — confundir electrón-volt con volt.
\end{itemize}
El botón «Practicar» del portal usa estos códigos para llevarte a la escena exacta de EquipotentialLab donde se repite el error.

\esection{Bibliografía y conexión}
Young, H.~D. y Freedman, R.~A. (2009). \emph{Física universitaria con física moderna}, Vol.~2 (12.\textsuperscript{a} ed.). Pearson. Capítulo~23 ``Potencial eléctrico'', secciones 23.1--23.5, pp.~780--805, ejemplos 23.1--23.5.

\vspace{0.6em}
\begin{center}
  \fcolorbox{e10line}{e10accentsoft}{\parbox{0.92\linewidth}{\centering
    \textbf{Conexión curricular:} $\ $ viene de M02 (Campo eléctrico) $\ \rightarrow\ $ \textbf{M03 (Potencial)} $\ \rightarrow\ $ lleva a M04 (Distribuciones continuas) y M06 (Capacitancia).}}

  \small\textcolor{e10muted}{Simulador: EquipotentialLab · Taller: \texttt{recursos/Cap03\_Taller.pdf} · © 2026 Dudbil Olvasada Pabón Riaño · 10.Easy}
\end{center}

\end{document}
